% ch-a6-06.tex —— 第 6 篇（Feng–Harris）续：PDF 544–555 / 印刷 531–542
% 本片自 PDF 544 页首（印刷 531）起译；开头第一句是上一页中断的句子（§8.4.1 中古典点 x 的讨论）。

\ysrc{531}其中 $L/\Q_p$ 是某个有限扩张，$x$ 在其上有定义。点 $x$ 还依赖于一项次级数据：一个加细（refinement）$t_x$，亦即晶体 Frobenius 特征值的一个排序；下文默认假定已给定该数据。表示 $\rho_\pi$ 并不依赖于 $t_x$ 的选取。

接着我们按 §8.1 的做法定义对偶伴随 Selmer 群（dual adjoint Selmer group）$H^1_f(\Q,M_\pi(1))=H^1_f(\Q,\operatorname{Ad}\circ\rho_\pi(1))$。注意这是一个 $\Q_p$-向量空间；与 [GV18,Ven19,HV19] 的构造不同，[HT17] 中的构造完全在特征零中展开。

[HT17] 的主要结果是：至少在有利条件下，并假定下文定理 8.4.1 的陈述中回顾的那个猜想成立，$H^1_f(\Q,M_\pi(1))$ 控制着 $\mathcal E$ 在 $x$ 处相对于态射 $\mathrm{pr}$ 的无穷小几何。把这一点与 [Han17] 中构造的泛系数谱序列结合起来，作者们重新得到了 Galatius–Venkatesh 定理 8.2.1 的类比，但不引入导出形变环。

[HT17] 最终结果所需的主要假设，是同构
\begin{equation}\tag{8.4.2}\label{eq:a6-06-1}
R_x\isom\mathbb T_x
\end{equation}
的一个版本，该同构来自 Taylor–Wiles 方法。这里 $R_x$ 是用 Robba 环上 $(\varphi,\Gamma)$-模的理论定义的非导出（underived）形变环。特别地，$R_x$ 完全由特征零的方法定义，并且（pro-）表示一个函子，该函子定义在 Artin 局部 $L$-代数上，$L$ 是 $\rho_\pi$ 的 $p$-进系数域（见 (8.4.1)）。更确切地说，假定 $\rho_\pi$ 绝对不可约；此时 $R_x$（pro-）表示 $\rho_\pi$ 连同加细 $t_x$ 的三角化（trianguline）形变。三角化性\footnote{英语读音与 fine、wine、turpentine 押韵。} [Ber11] 由 $\rho_\pi$ 限制到 $p$ 处分解群所得的 $(\varphi,\Gamma)$-模决定；对我们来说，相关的事实是：附着在特征簇上各点的伽罗瓦表示（包括 $\rho_\pi$ 本身）必然具有这一性质。

现在我们假定 $\pi$ 对 $H^*(S_{K_0},\tilde V)$ 有贡献，记号同 §7.4，其中 $G=\mathrm{GL}(n)_{\Q}$。我们略微改动记号：以 $\mathfrak m_\pi$ 记 $\mathbb T_{K_0}\otimes\Q_p$ 中对应于 $\pi$ 的极大理想，以 $H^*(S_{K_0},\tilde V)_{\mathfrak m_\pi}$ 记相应的局部化。令 $q_0=q_0(\mathrm{GL}(n)_{\Q})$、$\ell_0=\ell_0(\mathrm{GL}(n)_{\Q})$。下述定理总结了 [HT17] 定理 4.9 与 4.13 的主要结果。

\begin{thm}{8.4.1}\label{thm:a6-06-1}
设 $x\in\XX$ 是如上所述的古典点，$w=p(x)$，而 $\lambda:\mathbb T\to L$ 是相应的 Hecke 代数特征。假设 $\lambda$ 数值非临界\footnote{这是一个正则性条件，参见 [HT17, Theorem 4.7]。}。则
\begin{enumerate}[(1)]
\item $\dim\mathbb T_x\ge\dim\Lambda_w-\ell_0$；
\item 若 $\dim\mathbb T_x=\dim\Lambda_w-\ell_0$，令 $V_x=\ker(\Lambda\to\mathbb T_x)\otimes_{\Lambda_w}L$，其中 $\Lambda_w\to L$ 是增广映射。则 $H^*(S_{K_0},\tilde V)_{\mathfrak m_\pi}$ 作为外代数 $\wedge^*_LV_x$ 上的分次模自由、秩为 $1$，并由它在最高次数 $q_0+\ell_0$ 处的项生成。
\item 假设 $\rho_\pi$ 绝对不可约、在 $p$ 处晶体，具有预期的 Hodge–Tate 权；并且对每个素数 $\ell$（含 $p$），其相伴随的 Weil–Deligne 表示在局部 Langlands 对应下对应于 $\pi_\ell$。再假设同构 (8.4.2) 成立。
\end{enumerate}

\ysrc{532}则 $V_x$ 可与 $H^1_f(\Q,\operatorname{Ad}\rho_\pi(1))=\mathbf v_{\pi,p}[1/p]$ 等同起来，并且 $H^*(S_{K_0},\tilde V)_{\mathfrak m_\pi}$ 作为外代数 $\wedge^*_L\mathbf v_{\pi,p}$ 上的分次模自由、秩为 $1$，由它在最高次数 $q_0+\ell_0$ 处的项生成。\tn{底本此处及以下用天城文字母（梵文 va，转写作 v）排出 $\mathbf v_{\pi,p}$ 这一记号，中译本改排为粗体 $\mathbf v$。}
\end{thm}

\begin{rem}{8.4.2}\label{thm:a6-06-2}
自然可以假定，$\wedge^*_L\mathbf v_{\pi,p}$ 在 $H^*(S_{K_0},\tilde V)_{\mathfrak m_\pi}$ 上的作用在对偶意义下与定理 8.2.1 中定义的作用一致\tn{底本此处作 $H^*(S_{K_0},\tilde V))_{\mathfrak m_\pi}$，多一个右括号（原书 p.532）。}。然而，$V_x$ 与 $\mathbf v_{\pi,p}[1/p]$ 的等同依赖于形变环 $R_x$，而 $R_x$ 又依赖于加细 $t_x$ 的选取。Hansen 猜想这一作用不依赖于该额外选取，但即便这一猜想也仍未解决。\footnote{感谢 David Hansen 指出这一点。}
\end{rem}

\subsubsection{Khare–Ronchetti 构造}

本节的构造基于论文 [KR23b]。他们的导出菱形代数（derived diamond algebra）虽纯局部地定义，但这一构造与整体上同调的普通部分上的作用紧密相连。导出菱形代数的表示论如何纳入假想中的 $p$-进局部 Langlands 纲领，我们尚不清楚。

设 $p>2$ 是素数，$G$ 是有限扩张 $F/\Q$ 上的群 $\mathrm{PGL}(n)$。对任意开紧子群 $K\subset G$ 与任意环 $A$，以 $H(G,K;A)$ 记 $G$ 上 $K$-双不变函数的（非导出）Hecke 代数，系数在 $A$ 中。

设 $B=T\cdot U$ 是上三角 Borel 子群，$U$ 是它的幂幺根，$T$ 是对角极大环面。我们只关注 $\mathrm{PGL}(n)$，以便使用显式的矩阵群，并利用已知的、附着于尖点上同调类的伽罗瓦参数的构造。[KR23b] 处理了一般情形，读者应能自行作适当调整。

设 $U^-$ 是对立的幂幺（下三角）子群。设 $v$ 是 $F$ 中整除 $p$ 的素位，整数环 $\OO=\OO_v$，单值化元 $\varpi=\varpi_v$。对 $1\le b\le c$，令
\begin{equation}\tag{8.4.3}\label{eq:a6-06-2}
I_{b,c,v}=U^-(\varpi^c\OO)\cdot T(1+\varpi^b\OO)\cdot U(\OO);
\end{equation}
也就是说，它是同余子群
\[
I_{b,c,v}=\begin{pmatrix}1+\varpi^b*&*\\ \varpi^c*&1+\varpi^b*\end{pmatrix}\subset\mathrm{GL}(2,\OO)
\]
从 $\mathrm{GL}_2$ 到 $G$ 的推广。我们令 $I_{b,c}=\prod_{v\mid p}I_{b,c,v}$。

设 $K=K_p\cdot K^p\subset G(\mathbf A_f)$，其中 $K^p$ 是 $p$ 之外固定的水平子群，$K_p$ 取遍 (8.4.3) 中引入的子群 $\prod_{v\mid p}I_{b,c,v}$。我们考虑维数为 $d$ 的局部对称空间的上同调
\[
Y_{b,c}=G(\Q)\backslash G(\mathbf A_f)/K_\infty\cdot K^p\cdot I_{b,c},
\]
其中 $K_\infty\subset G(\R)$ 是极大紧子群，$K^p$ 是 $p$ 之外固定的水平子群，于是 $K^p\cdot I_{b,c}\subset G(\mathbf A_f)$ 是开紧的。我们假定 $K^p$ 在一个有限素数集 $S$ 之外是超特殊极大子群。

$Y_{b,c}$ 的上同调是非分歧 Hecke 代数
\[
H^{\mathrm{un}}=\prod_{q\notin S,\ q\neq p}H(G_q,K_q)
\]
上的模。它也是 $p$ 处某个导出 Hecke 代数上的模；下面我们在 Hida 理论的框架中定义该代数。

\ysrc{533}仿照 Hida，我们定义一个算子 $U_p=\prod_{\alpha\in\Sigma}U_\alpha$，它在每个 $Y(K(b,c))$ 上作为对应作用，并与自然映射 $Y(K(b,c))\to Y(K(b',c'))$（$b'<b$，$c'<c$）相容。（真正的算子应记作 $U_p(b,c)$，但我们不强调这一点。）粗略地说，当 $G=\mathrm{PGL}(n)$ 时有
\[
U_p=\prod_{j=1}^{n-1}U_{j,p},\qquad U_{j,p}=U_{\alpha_j}=I_{b,c}\prod_{v\mid p}\operatorname{diag}(\varpi_v\mathrm{Id}_j,\mathrm{Id}_{n-j})I_{b,c}.
\]

\begin{defn}{8.4.3}\label{thm:a6-06-3}
普通部分 $H^i(Y_{b,c},\Z/p^n\Z)_{\mathrm{ord}}\subset H^i(Y_{b,c},\Z/p^n\Z)$ 是 $U_p$ 算子可逆地作用的子模（并且是直和项）。
\end{defn}

\begin{fact}{8.4.4}\label{thm:a6-06-4}
对一切 $c\ge b\ge1$、$n\ge1$，拉回
\[
H^i(Y_{b,b},\Z/p^n\Z)_{\mathrm{ord}}\to H^i(Y_{b,c},\Z/p^n\Z)_{\mathrm{ord}}
\]
是同构。
\end{fact}

\begin{proof}
[KR23b] 中的证明是作用在奇异上链上的双陪集计算。
\end{proof}

\begin{rem}{8.4.5}\label{thm:a6-06-5}
若事实 8.4.4 能推广到超收敛上同调，则它必然需要以特征零的方式表述。
\end{rem}

群 $I_{1,c}/I_{b,c}$ 同构于阿贝尔 $p$-群
\[
T_b:=\prod_{v\mid p}T(1+\varpi_v\OO_v)/T(1+\varpi_v^b\OO_v).
\]
由于 $I_{b,c}$ 在 $I_{1,c}$ 中正规，元素 $\alpha\in I_{1,c}$ 定义菱形算子
\[
\langle\alpha\rangle=I_{b,c}\alpha I_{b,c}\in H_{\Z/p^m\Z}(G,I_{b,c}),
\]
对一切 $b,c$ 皆如此。

$\langle\alpha\rangle$ 在 $H^i(Y_{b,c},\Z/p^n\Z)$ 上的作用与映射 $Y_{b,c}\to Y_{b',c'}$（$b'<b$，$c'<c$）交换。令 $\Lambda_b=\Z_p[[T_b]]$，$\Lambda=\varprojlim_b\Lambda_b$，即 $T(\Z_p)$ 的 Iwasawa 代数。对任意有限扩张 $F/\Q_p$，我们还使用 $\Lambda_F=\Lambda\otimes_{\Z_p}F$。

$T_{b,c}$-覆盖 $Y_{c,c}/Y_{b,c}$ 对应于到分类空间的一个态射\tn{底本此处作 $T_{b,c}$，但上文只定义了 $T_b$（原书 p.533）。}
\[
\pi_{b,c}:Y_{b,c}\to BT_{b,c}.
\]
于是有环同态
\[
\pi_{b,c}^*:H^*(BT_{b,c},\Z/p^m\Z)=H^*(T_{b,c},\Z/p^m\Z)\to H^*(Y_{b,c},\Z/p^m\Z),
\]
从而有 $H^*(T_{b,c},\Z/p^m\Z)$ 在 $H^*(Y_{b,c},\Z/p^m\Z)$ 上的作用，经由导出菱形算子实现。由于 $I_{b,c}^{ab}=T_{b,c}$，有同构
\[
\iota:H^1(T_{b,c},\Z/p^m\Z)\isom H^1(I_{b,c},\Z/p^m\Z)\subset\mathcal H_{\Z/p^m\Z}(G_p,I_{b,c}).
\]

\begin{lem}{8.4.6}（[KR23b]）\label{thm:a6-06-6}
$\phi\in H^1(T_{b,c},\Z/p^m\Z)$ 的导出菱形作用，与 $\iota(\phi)$ 的导出 Hecke 作用（按事实 6.2.1 得到）一致。特别地，当 $b=1$ 时，这一作用与 $U_{j,p}$ 算子交换。
\end{lem}

关于交换性的断言来自一个显式的双陪集计算。它蕴含（当 $Y_{b,c}=Y_{1,c}$ 时）

\begin{cor}{8.4.7}\label{thm:a6-06-7}
导出的次数 $1$ 菱形算子作用在 $H^*(Y_{1,c},\Z/p^m\Z)_{\mathrm{ord}}$ 上，并保持所有 $U_{j,p}$-特征空间。
\end{cor}

\ysrc{534}另一项计算表明：

\begin{prop}{8.4.8}\label{thm:a6-06-8}
当 $c$ 变动时，$H^1(T_{1,c},\Z/p^m\Z)$ 在 $H^*(Y_{1,1},\Z/p^m\Z)_{\mathrm{ord}}$ 上的作用彼此相容，因而在极限中给出 $H^1(T(1+p\Z_p),\Z/p^m\Z)$ 在 $H^*(Y_{1,1},\Z/p^m\Z)_{\mathrm{ord}}$ 上的作用\tn{底本此处作 $H^1(T(1+p\Z_p)),\Z/p^m\Z)$，多一个右括号（原书 p.534）。}。

令 $m\to\infty$，我们便得到 $H^*(T(1+p\Z_p),\Z_p)$ 在 $H^*(Y_{1,1},\Z_p)_{\mathrm{ord}}$ 上的分次作用。
\end{prop}

我们回到 §7.2 的设定。设 $\pi$ 是 $G$ 的尖点上同调表示，如该节所述。于是由推论 7.2.5，
\[
H^{q_0(G)+\bullet}_0(Y_{1,1},\overline\Q_p)[\pi_f]
\]
是（$p$-进的）$\wedge^\bullet(\mathfrak a^*)$ 上的自由微分分次模，秩为 $0$ 或 $1$；这里我们用到了 $\mathrm{PGL}(n)$ 的重数一性质以及 Whittaker 模型的存在性。

下述定理是 [KR23b] 主要结果的一个不精确版本。设 $E$ 是 $\pi_f$ 的适当的 $p$-进系数域。

\begin{thm}{8.4.9}（[KR23b]）\label{thm:a6-06-9}
在某些有利情形下，$H^*(Y_{1,1},E)_\pi$ 由导出菱形算子 $H^*(T(1+p\Z_p),E)$ 的作用在 $H^{q_0}(Y_{1,1},E)_\pi$ 上生成。

此外，存在从 $H^1(T(1+p\Z_p),E)$ 到附着于 $\pi$ 的 Selmer 群 $\mathbf v_{\pi,p}[1/p]^\vee$ 的典范映射。并且在更有利的情形下，$H^*(T(1+p\Z_p),E)$ 在 $H^*(Y_{1,1},E)_\pi$ 上的作用会通过这个 Selmer 群分解。
\end{thm}

定理所暗示的「有利情形」包括同构 (8.4.2)，换言之即非阿贝尔 Leopoldt 猜想。与定理 8.2.3 一样，定理 8.4.9 与定理 8.4.1 的作用应当互为对偶。

\section{开放问题选}

我们讨论导出代数几何与 Langlands 对应的交汇处的一些开放问题。

\subsection{导出的局部形变条件}

本问题的背景是 §3 与 §5。在 §3 中我们讨论了若干模空间的例子，我们知道如何把它们升级为导出模空间。然而还有许多有趣的例子随具体问题一同出现，对它们我们不知道如何构造任何良好的导出增强，作者们也不清楚这样的增强是否应当存在。带局部条件的伽罗瓦形变函子这类例子特别有趣，因为它们大概会对自守性提升产生影响。

实践中，在实施模性提升论证时，人们希望对伽罗瓦形变函子施加更一般的局部条件，例如本论文集中 [CS23] 的文章所述。做法是为带局部条件的局部域构造伽罗瓦形变函子，结果发现它们往往不是 LCI；「$\ell=p$」的例子见 [CG18, Theorem 3.19] 与 [Sno18]，$\ell\neq p$ 的例子见 [BKM21]。这一现象值得理解，因为 Kisin 对 Taylor–Wiles 方法的改型把自守性提升「归约」为对局部形变环的良好控制问题。

\ysrc{535}我们不知道这些局部形变函子是否真的应当受隐藏光滑性哲学支配，因为它们通常不是「诚实」的模函子：其模空间描述并不「具体」地已知。相反，它们通常从「诚实」的模空间出发，经由概形论像或 Zariski 闭包之类的过程构造出来，而这些过程没有良好的导出版本。

例如，$p$ 处的局部形变函子常常由下述程序构造（据我们所知，这一做法首见于 [CDT99]）：用 $p$-进 Hodge 理论在无限制形变函子的 $\overline\Q_p$-点中定义一个子集，再取其 Zariski 闭包。（要点在于：$p$-进 Hodge 理论在系数为特征零的表示这一框架中被理解得最为一般。）于是特征零点具有某种模论描述，即用带 $p$-进 Hodge 理论条件的理性伽罗瓦表示来描述；但先验地不清楚如何具体地用伽罗瓦表示解释 $\Z_p$-点或 $\F_p$-点。

有些情形下理论在整系数上表现良好，例如 Fontaine–Laffaille 范围。此时局部形变函子有完整的模论描述，甚至能看出它是光滑的。

一般而言，$p$ 处的局部形变函子由 Kisin [Kis08] 构造，但其构造同样涉及从一般纤维取概形论像或 Zariski 闭包之类的过程，我们不知道如何以导出方式施行这些过程。这就是我们不知道导出这些局部形变函子的自然途径的原因。不过过去几年，由于 Bhatt–Morrow–Scholze 与 Bhatt–Scholze 的工作，我们对整 $p$-进 Hodge 理论的理解有了重大进展。特别地，Bhatt–Scholze 在 [BS23] 中为晶体伽罗瓦表示中的格建立了新的解释，用棱镜 $F$-晶体（prismatic $F$-crystal）的语言表述。这本身不足以给出系数一般的晶体伽罗瓦表示的概念，但它是通向这样一个概念、进而通向晶体形变函子的模论描述的有希望的一步；后者随后就可以导出。

然而我们要提醒：即使这一切都可行，目前也还不清楚这会带来什么应用。整体方法（例如 Taylor–Wiles 方法）的前提是：通过增加水平结构，整体形变环与自守形式的结构可以在适当意义上被「简化」。然而这种简化完全发生在局部到整体的方面，而若局部形变环是非经典的，就无法指望得到同类的简化。Iyengar–Khare–Manning 的工作 [IKM22] 并未直接处理这一顾虑，但它用导出代数推广了 Wiles 的数值判别法，是这一大方向上有希望的一步。

\subsection{微分分次 Hecke 代数}

本节讨论的背景是 §6.1。在 §6.1.5 中我们说明，局部导出 Hecke 代数 $\mathcal H(G(K);\Lambda)$ 的交换性仍属未知。我们猜想它一般成立，并且应当是微分分次 Hecke 代数 $\mathscr H(G(K);\Lambda)$ 的「高阶交换性」的推论。对微分分次代数，交换性是一种结构，而不是一种性质；事实上存在一族结构衡量一个代数「有多交换」，称为 $\mathbb E_n$-结构。结合（associative）微分分次代数带有 $\mathbb E_1$-结构，而 $\mathbb E_\infty$-结构是交换性的同伦相容类比。由构造可见 $\mathscr H(G(K);\Lambda)$ 有 $\mathbb E_1$-结构。若这一 $\mathbb E_1$-结构能提升为 $\mathbb E_2$-结构，则 $\mathcal H(G(K);\Lambda)$ 的交换性随之成立。我们相信情形如此，并进一步猜想 $\mathscr H(G(K);\Lambda)$ 上的 $\mathbb E_1$-结构可提升为 $\mathbb E_3$-结构。

\ysrc{536}我们来解释这一猜想的动机。Feng–Gaitsgory 的未发表工作 [FG]（有条件地）推出：若 $K$ 是函数域，则只要 $\Lambda$ 的特征 $\ell$ 相对于 $G$ 不太大，$\mathcal H(G(K);\Lambda)$ 就交换。[FG] 的策略是把 Frobenius 的范畴迹施用于模导出 Satake 等价；这一结果尚未见于文献，但 Arinkin–Bezrukavnikov 已宣布，且很可能可以直接从 Bezrukavnikov–Riche 的近作 [BR24] 推出。

这一策略并非只是把经典 Hecke 代数交换性的一个证明加以细化，而是从范畴等价出发。它应当还能在对偶群的语言中阐明导出 Hecke 代数的「含义」，但答案不易陈述，故我们不在此表述。这一策略的变体很可能对 $p$-进局部域 $K$ 也适用，不过要做到这一点，还需奠定相当多的基础工作。

现在说明为何这提示 $\mathscr H(G(K);\Lambda)$ 上可能存在 $\mathbb E_3$-结构。在上述策略中，微分分次 Hecke 代数作为导出 Satake 范畴的范畴迹出现。后者被期待具有 $\mathbb E_3$-结构，其中一个 $\mathbb E_1$-幺半运算来自卷积，另一个相容的 $\mathbb E_2$-幺半运算则与它作为代数曲线上「因子化范畴」的结构有关，后者是 $\mathbb E_2$-结构的代数类比。

\subsection{导出伽罗瓦形变环的作用}

本节讨论的背景是 §5。那里我们勾画了 [GV18] 的主要结果：构造导出伽罗瓦形变环 $\mathcal R^{\overline\rho,\mathrm{crys}}_{\OO_F[1/S]}$ 的同伦群在局部对称空间的同调群上的作用（该同调群在 Hecke 代数相应的极大理想处局部化）。

上述作用发生在经典代数层面：分次交换环作用在分次阿贝尔群上。我们期待这一作用可以细化为导出伽罗瓦形变环 $\mathcal R^{\overline\rho,\mathrm{crys}}_{\OO_F[1/S]}$ 在同调链的局部化上的作用（即动态交换环作用在动态阿贝尔群上），它在取同伦群之后诱导出前一作用。

这一作用的构造是一个引人入胜的挑战。回忆 Galatius–Venkatesh 的构造使用 Taylor–Wiles 拼接，其基础是在 $p$ 之外增加水平结构。在道德上，Galatius–Venkatesh 发现导出伽罗瓦形变环在拼接之后应当变成经典的，于是从经典作用推出导出作用；但这话字面上说不通，因为并不存在拼接后的导出伽罗瓦形变环。替代做法是：在 $p$ 处加上某种无穷水平结构，并希望相应的导出伽罗瓦形变环是经典的。确实，Khare–Ronchetti 的工作 [KR23b] 提示：若只考虑普通形变，则在沿普通（Hida）塔上升时应当会发生这种情况。另一方面，[KR23b] 也指出这一进路与 $p$-进超越性问题（例如 Leopoldt 猜想）密切相连，而众所周知这类问题极其困难，所以或许应当另寻角度。

\ysrc{537}尽管这一加细作用的构造看似同伦论问题，它却会对经典代数产生具体推论。事实上，借助谱 Hecke 代数 [Fen20] 的考虑，它实际会推出导出 Hecke 代数作用的 Venkatesh 互反律 [Ven19, Theorem 8.5]。这正是我们想要的，因为在所引定理目前的表述中，并没有有效的办法判定它对哪些素数 $q$ 适用。

\subsection{Hitchin 叠的迹猜想}

本节讨论的背景是 §4。下述问题摘自 [FK24, §6.4]。设
\begin{equation}\tag{9.4.1}\label{eq:a6-06-3}
\begin{tikzcd}[column sep=large]
\mathcal M & \mathcal C\arrow[l,"c_0"']\arrow[r,"c_1"] & \mathcal M
\end{tikzcd}
\end{equation}
是 $\F_q$ 上（导出）叠的对应。用下述导出笛卡儿方
\begin{equation}\tag{9.4.2}\label{eq:a6-06-4}
\begin{tikzcd}[column sep=large, row sep=large]
\mathrm{Sht}_{\mathcal M}\arrow[r,"\Delta'"]\arrow[d,"c'"'] & \mathcal C\arrow[d,"{(\Frob\circ c_0,c_1)}"]\\
\mathcal M\arrow[r,"\Delta"'] & \mathcal M\times\mathcal M
\end{tikzcd}
\end{equation}
来定义 $\mathrm{Sht}_{\mathcal M}$。

[FYZ23, Lemma 4.2.1] 证明了：

\begin{lem}{9.4.1}\label{thm:a6-06-10}
切复形 $\mathbf T_{\mathrm{Sht}_{\mathcal M}/\F_q}$ 是 $\mathbf T_{c_1}$ 的限制。特别地，若 $c_1$ 拟光滑，则 $\mathrm{Sht}_{\mathcal M}$（在 $\Spec\F_q$ 上）拟光滑，维数等于 $d(c_1)$（$c_1$ 的相对虚拟维数）。
\end{lem}

我们将对母题层（motivic sheaf）表述迹猜想。母题层的相关定义见 [FK24, §3]；不熟悉母题层理论的读者可以在下文的讨论中把 $\mathbf Q$ 换成常值层 $\mathbf Q_\ell$，而不失要点。

对母题层 $\mathcal K$ 与 $i\in\Z$，记 $\mathcal K\langle i\rangle:=\mathcal K[2i](i)$。对（导出）叠的态射 $\mathcal X\to\mathcal Y$，记 $\operatorname{CH}_i(\mathcal X/\mathcal Y):=H^{-2i}(\mathcal X,f^!\mathbf Q_{\mathcal Y}(-i))$。当 $f\colon\mathcal X\to\Spec k$ 是到域的结构态射时，简记 $\operatorname{CH}_i(\mathcal X):=\operatorname{CH}_i(\mathcal X/\Spec k)$。若 $f$ 拟光滑，则记 $d(f)=\chi(\mathbf T_f)$ 为其切复形的欧拉示性数，并称之为 $f$ 的「相对（虚拟）维数」。由 [Kha19]，此时存在相对基本类 $[f]\in\operatorname{CH}_{d(f)}(\mathcal X/\mathcal Y)$。

假设在上述设定中 $c_1$ 拟光滑，于是 $[c_1]\in\operatorname{CH}_{d(c_1)}(\mathcal C/\mathcal M)$ 存在。则由引理 9.4.1，$\mathrm{Sht}_{\mathcal M}$（在 $\Spec\F_q$ 上）拟光滑，故导出基本类 $[\mathrm{Sht}_{\mathcal M}]\in\operatorname{CH}_{d(c_1)}(\mathrm{Sht}_{\mathcal M})$ 存在。

另一方面，把 $[c_1]$ 视为映射 $c_1^*\mathbf Q_{\mathcal M}\to c_1^!\mathbf Q_{\mathcal M}\langle-d(c_1)\rangle$，我们有一列映射
\[
(\Frob\circ c_0)^*\mathbf Q_{\mathcal M}=\mathbf Q_{\mathcal C}=c_1^*\mathbf Q_{\mathcal M}\xrightarrow{[c_1]}c_1^!\mathbf Q_{\mathcal M}\langle-d(c_1)\rangle,
\]
其复合记为 $\mathfrak c_{\mathcal M}$。这是 $\mathbf Q_{\mathcal M}$ 与 $\mathbf Q_{\mathcal M}\langle-d(c_1)\rangle$ 之间的一个上同调对应。我们将定义另一个类 $\operatorname{Tr}(\mathfrak c_{\mathcal M})\in\operatorname{CH}_{d(c_1)}(\mathrm{Sht}_{\mathcal M})$。

\ysrc{538}考虑 (9.4.2)，并把右侧竖直映射简记为 $c$。[FK24, §3.8] 中说明
\begin{equation}\tag{9.4.3}\label{eq:a6-06-5}
\begin{aligned}
\mathcal R\mathcal H\mathrm{om}\big((\Frob\circ c_0)^*\mathbf Q_{\mathcal M},c_1^!\mathbf Q_{\mathcal M}\big)
&\cong c^!\mathcal R\mathcal H\mathrm{om}\big(\mathrm{pr}_0^*\mathbf Q_{\mathcal M},\mathrm{pr}_1^!\mathbf Q_{\mathcal M}\big)\\
&\cong c^!\big(\mathbb D(\mathbf Q_{\mathcal M})\boxtimes\mathbf Q_{\mathcal M}\big),
\end{aligned}
\end{equation}
其中 $\mathbb D$ 是 Verdier 对偶函子。求值映射 $\mathbb D(\mathbf Q_{\mathcal M})\otimes\mathbf Q_{\mathcal M}\langle-d(c_1)\rangle\to\mathbb D_{\mathcal M}\langle-d(c_1)\rangle$ 伴随于一个映射 $\mathbb D(\mathbf Q_{\mathcal M})\boxtimes\mathbf Q_{\mathcal M}\to\Delta_*\mathbb D_{\mathcal M}\langle-d(c_1)\rangle$，其中 $\mathbb D_{\mathcal M}$ 是 $\mathcal M$ 的对偶复形。把它与 (9.4.3) 复合，得到映射
\begin{equation}\tag{9.4.4}\label{eq:a6-06-6}
\mathcal R\mathcal H\mathrm{om}\big((\Frob\circ c_0)^*\mathbf Q_{\mathcal M},c_1^!\mathbf Q_{\mathcal M}\big)\to c^!\Delta_*\mathbb D_{\mathcal M}\langle-d(c_1)\rangle.
\end{equation}
最后，利用正常基变换，有同构
\begin{equation}\tag{9.4.5}\label{eq:a6-06-7}
\begin{aligned}
c^!\Delta_*\mathbb D_{\mathcal M}\langle-d(c_1)\rangle
&\cong\Delta'_*(c')^!\mathbb D_{\mathcal M}\langle-d(c_1)\rangle\\
&\cong\Delta'_*\mathbb D_{\mathrm{Sht}_{\mathcal M}}\langle-d(c_1)\rangle.
\end{aligned}
\end{equation}
我们可以把 $\mathfrak c_{\mathcal M}$ 视为 $\mathcal R\mathcal H\mathrm{om}(c_0^*\mathbf Q_{\mathcal M},c_1^!\mathbf Q_{\mathcal M}\langle-d(c_1)\rangle)$ 的一个整体截面。于是
\[
\operatorname{Tr}(\mathfrak c_{\mathcal M})\in H^0\big(\mathcal C,\Delta'_*\mathbb D_{\mathrm{Sht}_{\mathcal M}}\langle-d(c_1)\rangle\big)\cong\operatorname{CH}_{d(c_1)}(\mathrm{Sht}_{\mathcal M})
\]
定义为它在 (9.4.4) 与 (9.4.5) 的复合下的像。

至此，我们在 $\operatorname{CH}_{d(c_1)}(\mathrm{Sht}_{\mathcal M})$ 中有了两个自然的类：一个是导出基本类 $[\mathrm{Sht}_{\mathcal M}]$，另一个是 $\operatorname{Tr}(\mathfrak c_{\mathcal M})$。于是自然会问：
\[
\operatorname{Tr}(\mathfrak c_{\mathcal M})=[\mathrm{Sht}_{\mathcal M}]\in\operatorname{CH}_{d(c_1)}(\mathrm{Sht}_{\mathcal M})
\]
是否成立。我们没有证据相信它在所述的全部一般性下成立。然而，当 $\mathcal M$ 是例 3.2.14 中考虑的那类导出 Hitchin 叠时，我们觉得这一等式多半是模性猜想 [FYZ21, Conjecture 4.12] 成立所必需的，因此我们猜想该等式在此情形下成立。

\begin{conj}{9.4.2}（迹猜想）\label{thm:a6-06-11}
设 $\mathcal M$ 是 [FYZ21, §5] 中的导出 Hitchin 叠。则
\[
\operatorname{Tr}(\mathfrak c_{\mathcal M})=[\mathrm{Sht}_{\mathcal M}]\in\operatorname{CH}_{d(c_1)}(\mathrm{Sht}_{\mathcal M}).
\]
\end{conj}

这样的 $\mathcal M$ 具有一些特殊性质，这些性质或许对证明是必要的；例如它们都是拟光滑的\footnote{注意，这一拟光滑性条件与迹猜想的表述无关。}。我们期待迹猜想会成为模性猜想 4.5.1 的证明中的重要一环。

\section*{附录 A.\quad 单纯交换环速成}

本附录是关于单纯交换环的一份粗略入门，最终引向余切复形。其目的是为那些觉得正文的「黑箱」进路过于含糊的读者提供更多根基，并为有兴趣深入这一主题的读者指明参考文献。简短的介绍我们喜欢 [Qui, Qui70, Iye07]；较长的教科书有 [May67] 与 [GJ09]。

特别地，我们认为，即使读者的最终目标是在动态环与 $\infty$-范畴的语言中工作，了解单纯交换环与模型范畴也是有益的。

我们的处理远非全面。事实上，我们把写作的粗略性视为它的主要（也许是唯一）特色。我们触及的每个主题在文献中都已记录充分，但单纯同伦论是一个如此庞大且技术性的主题，确有迷失于细节的风险。为与正文的精神一致，我们宁愿只勾画该主题的骨架，满足于口号与直观，而把细节留给我们所偏爱的参考文献。

\ysrc{539}\subsection*{A.1.\quad 为何用「单纯」？}

我们将把「单纯集」当作拓扑空间的一种组合模型来介绍，把「单纯交换环」当作拓扑交换环的一种组合模型来介绍。

为什么说「单纯」而不说「拓扑」？一个答案是：单纯对象的组合本性使它们在技术层面表现更好——它排除了「病态」的拓扑空间。我们没有好的哲学答案；读者可以参阅 [Matb] 中关于这一问题的更多思考。

此外，事实证明「单纯」这个形容词足够灵活，可以施用于任何范畴中的任何对象，并在许多有趣的例子中提供一种做抽象同伦论的方法。也就是说，若 $\mathsf C$ 是某个「widget」的范畴（widget 是占位词，代指任意一类对象），我们就能定义「单纯 widget」的范畴 $\mathsf{sC}$。

\subsection*{A.2.\quad 单纯范畴}

单纯范畴（simplex category）$\Delta$ 的对象是集合 $[n]:=\{0,1,\dots,n\}$（$n\ge0$），而 $\Hom_\Delta([n],[m])$ 由所有不反转序的集合映射 $f\colon[n]\to[m]$ 组成，即若 $i\le j$，则 $f(i)\le f(j)$。

我们将为 $[n]$ 与 $[n+1]$ 之间的生成态射命名。对 $0\le i\le n+1$，余面映射
\[
\delta_i\colon[n]\to[n+1]
\]
是「跳过 $i$」的态射，更正式地
\[
\delta_i(j)=\begin{cases}j & j<i,\\ j+1 & j\ge i.\end{cases}
\]
对 $0\le i\le n$，余退化映射
\[
\sigma_i\colon[n+1]\to[n]
\]
是「在 $i$ 处加倍」的态射，更正式地
\[
\sigma_i(j)=\begin{cases}j & j\le i,\\ j-1 & j\ge i+1.\end{cases}
\]

\begin{rem}{A.2.1}\label{thm:a6-06-12}
这些态射「生成」的意思是：$\Delta$ 中任何态射都能写成若干余退化映射的复合，再接若干余面映射的复合。
\end{rem}

\subsubsection*{A.2.1.\quad 单纯集}

单纯集是函子 $\underline X\colon\Delta^{\mathrm{op}}\to\mathsf{Sets}$。具体地说，它由一族集合 $X_n:=\underline X([n])$（$n\ge0$）以及由 $\Delta^{\mathrm{op}}$ 中态射索引的它们之间的映射给出。特别地，有面映射
\[
d_i:=\underline X(\delta_i)\colon X_{n+1}\to X_n\quad(0\le i\le n+1),
\]
与退化映射
\[
s_i:=\underline X(\sigma_i)\colon X_n\to X_{n+1}\quad(0\le i\le n).
\]

\ysrc{540}由于注 A.2.1，要确定一个单纯集，只需给出 $\{X_n\}_{n\ge0}$ 以及每个 $n$ 的 $d_i$ 与 $s_i$，使它们满足「单纯恒等式」[GJ09, I.1]：
\begin{itemize}
\item $d_id_j=d_{j-1}d_i$，$i<j$；
\item $d_js_j=d_{j+1}s_j=\mathrm{Id}$；
\item $d_is_j=s_{j-1}d_i$，$i<j$；
\item $d_is_j=s_jd_{i-1}$，$i>j+1$；
\item $s_is_j=s_{j+1}s_i$，$i\le j$。
\end{itemize}

\begin{exam}{A.2.2}\label{thm:a6-06-13}
$\Delta$ 的每个元素 $[n]$ 都诱导 $\Delta^{\mathrm{op}}$ 上的一个可表函子，即 $\Hom_\Delta(-,[n])$，从而诱导一个单纯集，我们称之为 $\Delta^n$。直观上，我们把它看作对应于标准 $n$-单形的拓扑空间，记作
\[
|\Delta^n|:=\{(x_0,\dots,x_n)\in\mathbb R^{n+1}\colon0\le x_i\le1,\ \textstyle\sum x_i=1\}.
\]
\end{exam}

\begin{exam}{A.2.3}\label{thm:a6-06-14}
单纯集的定义是对下述结构的公理化：取拓扑空间 $X$，令 $\underline X[n]$ 为从 $n$-单形 $|\Delta^n|$ 到 $X$ 的连续映射的集合，即 $X_n$ 是「$X$ 的 $n$-单形」的集合。面映射 $d_i$ 就是通常的边界映射。退化映射则给出把低维单形视为高维单形的「退化」实例的各种方式；它们在拓扑理论中并不重要，但在单纯理论中技术上很方便。这定义了奇异单形函子
\[
\Sing\colon\mathsf{Top}\to\mathsf{sSets}.
\]
\end{exam}

因此，对一般的单纯集 $\underline X$，我们把 $\underline X[n]$ 称为「$\underline X$ 的 $n$-单形」。

\subsubsection*{A.2.2.\quad 几何实现}

给定单纯集 $\underline X$，我们可以把 $\underline X$ 看作拼装单纯复形的配方，从而由它构造一个拓扑空间 $X:=|\underline X|$。用公式表示，几何实现由下式给出
\[
|\underline X|:=\frac{\coprod_n\underline X[n]\times|\Delta^n|}{(d_ix,u)\sim(x,\delta_iu),\ (s_ix,u)\sim(x,\sigma_iu)}.
\]
拓扑空间 $|\underline X|$ 是 CW 复形 [May67, §14]。

另一种简洁的表述如下：我们宣布 $|\Delta^n|$ 是标准 $n$-单形 $\{(x_0,\dots,x_n)\in\mathbb R^{n+1}\colon0\le x_i\le1,\sum x_i=1\}$，然后定义
\[
|\underline X|=\varinjlim_{\Delta^n\to\underline X}|\Delta^n|,
\]
其中指标范畴的对象是各 $n$ 对应的映射 $\Delta^n\to\underline X$，态射是尊重给定到 $\underline X$ 的映射的映射 $\Delta^m\to\Delta^n$。注意单纯集 $\Delta^n$ 的几何实现与标准 $n$-单形 $|\Delta^n|$ 一致，这说明了记号的合理性。

\begin{prop}{A.2.4}\label{thm:a6-06-15}
几何实现是奇异单形函子的左伴随：
\[
\Hom_{\mathsf{Top}}(|\underline X|,Y)\cong\Hom_{\mathsf{sSet}}(\underline X,\Sing Y).
\]
\end{prop}

\begin{rem}{A.2.5}\label{thm:a6-06-16}
在适当意义上，这两个伴随函子定义了范畴之间的一个「同伦等价」，而我们尚无工具精确地定义这一点。技术性的陈述是：这两个函子诱导模型范畴的 Quillen 等价，两侧都取通常的 Quillen 模型结构。精确表述与证明见 [May92, §16]。这就证实了单纯集提供拓扑空间的一个组合模型这一断言。
\end{rem}

\ysrc{541}\subsubsection*{A.2.3.\quad 单纯集的同伦群}

带基点的单纯集 $(\underline X,\star)$ 有同伦群 $\pi_i(\underline X,\star)$，它可以定义为几何实现 $|\underline X|$ 关于基点 $\star$ 的同伦群（通常意义下）。它也可以直接在单纯集的范畴中定义，基本做法是用球面的组合模型，即把单纯集 $\Delta^n$ 模去其边界；但这涉及用一个 Kan 复形「消解」$\underline X$ 的微妙之处（见注 A.2.5）。我们指出，对拓扑空间 $Y$，$\Sing Y$ 是 Kan 复形。要点在于：在「同伦论」范畴（例如单纯集的范畴）中，对象可能没有「足够多」的进出映射，因此必须用具有足够多映射的对象来「消解」。每个拓扑空间都是纤维性的，这就是为什么 $|\underline X|$ 不必类似地「消解」。

当我们讨论单纯（阿贝尔）群与单纯交换环时，总取 $\star$ 为对应于单位元的点，并在记号中省略基点。

\subsubsection*{A.2.4.\quad 内 Hom}

给定两个单纯集 $\underline X$ 与 $\underline Y$，我们构造一个单纯集 $\underline{\Map}(\underline X,\underline Y)$，其 $0$-单形的集合可等同于 $\Hom_{\mathsf{sSet}}(\underline X,\underline Y)$。直观上，我们是想在从 $\underline X$ 到 $\underline Y$ 的映射的集合上赋予拓扑结构。这将使单纯集的范畴成为一个单纯范畴，即充实于单纯集的范畴。

我们定义
\[
\underline{\Map}(\underline X,\underline Y)([n]):=\Hom_{\mathsf{sSet}}(\underline X\times\Delta^n,\underline Y),
\]
其中面映射与退化映射由 $\Delta^n$ 上的映射诱导。注意函子的乘积逐层构造。

\begin{rem}{A.2.6}\label{thm:a6-06-17}
单纯范畴是 $\infty$-范畴这一概念的一种可能模型（尽管不是 [Lur09] 中使用的主要模型）。直观上，这是一类充实于空间的范畴，即态射的集合被提升为空间。
\end{rem}

\subsection*{A.3.\quad 单纯 widget}

更一般地，若 $\mathsf C$ 是任意范畴，则称 $\mathsf C$ 的单纯对象为函子 $\Delta^{\mathrm{op}}\to\mathsf C$。若 $\mathsf C$ 是「widget」的范畴，则函子范畴 $\Hom(\Delta^{\mathrm{op}},\mathsf C)$ 将被称为「单纯 widget」的范畴，记作 $\mathsf{sC}$。

\begin{exam}{A.3.1}\label{thm:a6-06-18}
下述例子对我们特别重要。
\begin{itemize}
\item 单纯阿贝尔群是从 $\Delta^{\mathrm{op}}$ 到阿贝尔群范畴的函子。
\item 给定交换环 $R$（对我们而言，这总指有单位元的交换环），单纯 $R$-模是从 $\Delta^{\mathrm{op}}$ 到 $R$-模范畴的函子。
\item 单纯交换环是从 $\Delta^{\mathrm{op}}$ 到交换环范畴的函子。
\item 给定交换环 $R$，单纯 $R$-代数是从 $\Delta^{\mathrm{op}}$ 到 $R$-代数范畴的函子。
\end{itemize}
\end{exam}

\ysrc{542}这类对象构成的范畴在每种情形下都充实于 $\mathsf{sSet}$。为看出这一点，先注意：对单纯集 $\underline S$ 与单纯 widget $\underline X$，有单纯 widget
\[
(\underline X^{\otimes\underline S})[n]=\coprod_{S_n}X_n,
\]
其中右端的余积在 widget 的范畴中取，并带有显然的面映射与退化映射。然后定义
\[
\underline{\Hom}_{\mathsf{sC}}(\underline X,\underline Y)[n]:=\Hom_{\mathsf{sC}}(\underline X^{\otimes\Delta^n},\underline Y),
\]
同样带有显然的面映射与退化映射。

\begin{exam}{A.3.2}\label{thm:a6-06-19}
设 $\underline X$ 是单纯集。则有一个单纯阿贝尔群 $\mathbf Z\langle\underline X\rangle$，逐层取自由阿贝尔群得到：定义 $\mathbf Z\langle\underline X\rangle([n]):=\mathbf Z\langle\underline X([n])\rangle$，并让面映射与退化映射由 $\underline X$ 的相应映射诱导。

还有一个单纯 $\Z$-代数 $\mathbf Z[\underline X]$，逐层取自由（多项式）$\Z$-代数得到，面映射与退化映射由 $\underline X$ 的相应映射诱导。
\end{exam}

\begin{rem}{A.3.3}\label{thm:a6-06-20}
事实证明单纯群自动是 Kan 复形。因此，单纯群的同伦群可以按从单纯球面出发的映射来计算，这与 §A.2.3 中对一般单纯集的告诫情形相反。同样的结论对单纯阿贝尔群、单纯交换环等也成立，因为各类情形中模型范畴结构正是这样定义的。
\end{rem}

\subsection*{A.4.\quad Dold–Kan 对应}

回忆：若 $\mathsf C$ 是阿贝尔范畴，则可构造 $\mathsf C$ 中对象的链复形的范畴。这在同调代数中至关重要。事实证明，它与范畴 $\mathsf{sC}$ 密切相关。

\begin{thm}{A.4.1}（Dold–Kan）\label{thm:a6-06-21}
设 $R$ 是交换环。单纯 $R$-模的范畴等价于 $R$-模的非负分次链复形的范畴。
\end{thm}

我们可以把定理 A.4.1 看作在说明：「单纯」是把「链复形」推广到非阿贝尔范畴（那里「链复形」的概念不存在）。

我们指出定理 A.4.1 的等价所用到的函子，先从单纯 $R$-模到 $R$-模链复形的函子说起。证明可在 [GJ09, §III.2] 中找到；我们也喜欢 [Mata] 中的处理，并依循之。设 $\underline M$ 是单纯 $R$-模。先解释一个辅助构造，称为 $\underline M$ 的 Moore 复形。由 $\underline M$ 构造链复形 $M_*$：取 $M_n:=\underline M[n]$，微分 $\partial_n\colon M_n\to M_{n-1}$ 取面映射的交替和，
\[
\partial_n:=\sum_{i=0}^n(-1)^id_i.
\]

\begin{exam}{A.4.2}\label{thm:a6-06-22}
拓扑空间 $X$ 的奇异链复形，按定义就是单纯阿贝尔群 $\mathbf Z\langle\Sing X\rangle$ 的 Moore 复形。
\end{exam}

\begin{rem}{A.4.3}\label{thm:a6-06-23}
Moore 复形 $M_*$ 的同调群与单纯阿贝尔群 $\underline M$ 的同伦群一致。
\end{rem}
